%%%%%%%%%%%%%%%%%%%%%%%%%%%%%%%%%%%%%%%%%
% Short Sectioned Assignment LaTeX Template Version 1.0 (5/5/12)
% This template has been downloaded from: http://www.LaTeXTemplates.com
% Original author:  Frits Wenneker (http://www.howtotex.com)
% License: CC BY-NC-SA 3.0 (http://creativecommons.org/licenses/by-nc-sa/3.0/)
%%%%%%%%%%%%%%%%%%%%%%%%%%%%%%%%%%%%%%%%%
\UseRawInputEncoding
% \documentclass[paper=a4, fontsize=11pt]{scrartcl} % A4 paper and 11pt font size
\documentclass[11pt, a4paper]{book}
\usepackage[T1]{fontenc} % Use 8-bit encoding that has 256 glyphs
\usepackage[utf8]{inputenc}
\usepackage{fourier} % Use the Adobe Utopia font for the document - comment this line to return to the LaTeX default
\usepackage{listings} % para insertar código con formato similar al editor
\usepackage[spanish, es-tabla]{babel} % Selecciona el español para palabras introducidas automáticamente, p.ej. "septiembre" en la fecha y especifica que se use la palabra Tabla en vez de Cuadro
\usepackage{url} % ,href} %para incluir URLs e hipervínculos dentro del texto (aunque hay que instalar href)
\usepackage{graphics,graphicx, float} %para incluir imágenes y colocarlas
\usepackage[gen]{eurosym} %para incluir el símbolo del euro
\usepackage{cite} %para incluir citas del archivo <nombre>.bib
\usepackage{enumerate}
\usepackage{hyperref}
\usepackage{graphicx}
\usepackage{tabularx}
\usepackage{booktabs}

\usepackage[table,xcdraw]{xcolor}
\hypersetup{
	colorlinks=true,	% false: boxed links; true: colored links
	linkcolor=black,	% color of internal links
	urlcolor=cyan		% color of external links
}
\renewcommand{\familydefault}{\sfdefault}
\usepackage{fancyhdr} % Custom headers and footers
\pagestyle{fancyplain} % Makes all pages in the document conform to the custom headers and footers
\fancyhead[L]{} % Empty left header
\fancyhead[C]{} % Empty center header
\fancyhead[R]{Rubén Delgado Pareja} % My name
\fancyfoot[L]{} % Empty left footer
\fancyfoot[C]{} % Empty center footer
\fancyfoot[R]{\thepage} % Page numbering for right footer
%\renewcommand{\headrulewidth}{0pt} % Remove header underlines
\renewcommand{\footrulewidth}{0pt} % Remove footer underlines
\setlength{\headheight}{13.6pt} % Customize the height of the header

\usepackage{titlesec, blindtext, color}
\definecolor{gray75}{gray}{0.75}
\newcommand{\hsp}{\hspace{20pt}}
\titleformat{\chapter}[hang]{\Huge\bfseries}{\thechapter\hsp\textcolor{gray75}{|}\hsp}{0pt}{\Huge\bfseries}
\setcounter{secnumdepth}{4}
\usepackage[Lenny]{fncychap}


\begin{document}

	% Plantilla portada UGR
	\begin{titlepage}
\newlength{\centeroffset}
\setlength{\centeroffset}{-0.5\oddsidemargin}
\addtolength{\centeroffset}{0.5\evensidemargin}
\thispagestyle{empty}

\noindent\hspace*{\centeroffset}\begin{minipage}{\textwidth}

\centering
\includegraphics[width=0.9\textwidth]{logos/logo_ugr.jpg}\\[1.4cm]

\textsc{ \Large TRABAJO FIN DE GRADO\\[0.2cm]}
\textsc{ GRADO EN INGENIERIA INFORMATICA}\\[1cm]

{\Huge\bfseries Triage Inteligente de Consulta Médica \\}
\noindent\rule[-1ex]{\textwidth}{3pt}\\[3.5ex]
{\large\bfseries Sistema de citas priorizado por el riesgo del paciente }
\end{minipage}

\vspace{2.5cm}
\noindent\hspace*{\centeroffset}
\begin{minipage}{\textwidth}
\centering

\textbf{Autor}\\ {Rubén Delgado Pareja}\\[2.5ex]
\textbf{Director}\\ {Juan Julián Merelo Guervós}\\[2cm]
\includegraphics[width=0.3\textwidth]{logos/etsiit_logo.png}\\[0.1cm]
\textsc{Escuela Técnica Superior de Ingenierías Informática y de Telecomunicación}\\
\textsc{---}\\
Granada, Junio de 2022
\end{minipage}
\end{titlepage}


	% Plantilla prefacio UGR
	\thispagestyle{empty}

\begin{center}
{\large\bfseries Título \\ Subtítulo }\\
\end{center}
\begin{center}
Nombre Del Estudiante\\
\end{center}

%\vspace{0.7cm}

\vspace{0.5cm}
\noindent{\textbf{Palabras clave}: \textit{software libre}
\vspace{0.7cm}

\noindent{\textbf{Resumen}\\
	

\cleardoublepage

\begin{center}
	{\large\bfseries Same, but in English}\\
\end{center}
\begin{center}
	Student's name\\
\end{center}
\vspace{0.5cm}
\noindent{\textbf{Keywords}: \textit{open source}, \textit{floss}
\vspace{0.7cm}

\noindent{\textbf{Abstract}\\


\cleardoublepage

\thispagestyle{empty}

\noindent\rule[-1ex]{\textwidth}{2pt}\\[4.5ex]

D. \textbf{Tutora/e(s)}, Profesor(a) del ...

\vspace{0.5cm}

\textbf{Informo:}

\vspace{0.5cm}

Que el presente trabajo, titulado \textit{\textbf{Chief}},
ha sido realizado bajo mi supervisión por \textbf{Estudiante}, y autorizo la defensa de dicho trabajo ante el tribunal
que corresponda.

\vspace{0.5cm}

Y para que conste, expiden y firman el presente informe en Granada a Junio de 2018.

\vspace{1cm}

\textbf{El/la director(a)/es: }

\vspace{5cm}

\noindent \textbf{(nombre completo tutor/a/es)}

\chapter*{Agradecimientos}






	% Índice de contenidos
	\newpage
	\tableofcontents

	% Índice de imágenes y tablas
	\newpage
	\listoffigures

	% Si hay suficientes se incluirá dicho índice
	\listoftables
	\newpage

	% Introducción
	\chapter{Introducción}
Este proyecto es software libre, y está liberado con la licencia \cite{gplv3}.

Debido a la crisis económica del 2008 la sanidad española no ha parado de recibir recortes.
Puede observarse en este \href{https://www.consalud.es/politica/decada-recortes-2009-2018-efectos-infrafinanciacion-sanidad_87083_102.html}{artículo}
donde se estudia la financiación de la sanidad desde entonces. Esto provoca que disminuya el número de sanitarios en España, 
esta situación que se lleva dando desde hace más de una década produce una sobrecarga de trabajo para los sanitarios y un retardo 
en todas las listas de esperas, tanto en cirugias como en la atención primaria. 
Con la llegada de la pandemia y del covid-19 todo empeoró llegando a colapsar totalmente la atención primaria debido a la demanda
de citas previas y a la escasez de profesionales enfocados a este servicio. Además desde hace años los sanitarios se aglomeran
en los hospitales por mejores situaciones laborales y salarios. En las siguientes links podemos observar la situación descrita.
\begin{itemize}
    \item \url{https://elpais.com/sociedad/2021-02-08/la-pandemia-sume-a-la-atencion-primaria-en-una-saturacion-permanente.html} 
    \item \url{https://www.lavanguardia.com/vida/20211225/7952109/atencion-primaria-colapso-pacientes-medico-dias-espera.html}
    \item \url{https://www.lavozdegalicia.es/noticia/galicia/2022/01/16/dejando-atender-casos/0003_202201G16P2993.htm}

\section{Motivación}
El triage es una herramienta fundamental utilizada en urgenicas y emergencias para gestionar el riesgo clínico y manejar 
con seguridad los flujos de los pacientes en situaciones que las necesidades clínicas superan los recursos. 
Actualmente esta herramienta se utiliza con cinco niveles de prioridad que se asignan dependiendo de la urgencia 
del paciente, cuanto más urgente antes necesitan ser asistidos. 
La principal motivación de este proyecto es extrapolar esta herramienta y llevarla al campo de la informática, intentando liberar
la saturacion de las colas de espera, sin olvidar el riesgo de la salud de los pacientes.
A pesar de que el problema que se trata de abordar tenga grandes maginutes, la intención es desarrollar una herramienta útil para priorizar
los pacientes con mayores riesgos y que sirva de respaldo para los sanitarios centrados en la atención primaria
De esta forma, haciendo uso de los sistemas informáticos podamos resolver un problema de la sociedad

\section{Objetivos}

\subsection{Principal}

\textbf{Crear el triage inteligente de las consultas médicas}. El objetivo se centra en desarrollar un sistema informático que sirva
como apoyo a los sanitarios de atención primaria reduciendo, todo lo posible, la carga de trabajo de ellos, liberando de esta forma 
la saturación en las colas de esperas. Además al tratarse de un triage de consultas priorizaremos las personas con mayor riesgo 
aumentando la eficacia de la atención primaria. Gracias a este sistema podré aportar lo posible a la sanidad desde mi ámbito, la informática

El sistema informática consistirá en un triage de consultas médicas ; el triage en el ámbito de la sanidad es un proceso que permite gestionar
el flujo de los pacientes teniendo en cuenta el riesgo cuando la demanda y necesidades clínicas superan los recursos, es muy importante en urgencias clínicas
\href{https://scielo.isciii.es/scielo.php?script=sci_arttext&pid=S1137-66272010000200008}{Aquí} se puede encontrar más información

El sistema de triage conlleva crear un sistemas de colas de citas priorizada, las cualés se ordenarán mediante una heuristica que tendrá 
en cuenta el riesgo de cada usuario, también debe de estar integrada en una aplicación accesible para todos los usuarios y
asegurando la calidad del código

\subsection{Opcionales}
Los objetivos secundarios podrían ser:

\begin{itemize}
    \item Almacenar el historial médico de una consulta en la aplicación
    \item Modificar el riesgo de los usuario 

\subsection{Proyecto a futuro}
Los objetivos que se quedarán para el futuro del proyecto serán: 

\begin{itemize}
    \item Aplicar la aplicación para consultas especializadas 
    \item Construir un plugin para acceder a las recetas de los medicamentos de una consulta
    \item Crear consultas automatizadas para enfermedades crónicas


	% Descripción del problema y hasta donde se llega
	\UseRawInputEncoding
\chapter{Descripción del problema}
Desde el inicio de la pandemia, atender telemáticamente es fundamental para no saturar las consultas médicas, sin embargo si todo
se realiza de forma manual la saturación puede ser la misma. Mi proyecto intentará resolver este problema mediante un sistema de
triage el cual es capaz de filtrar solicitudes de atención médica con distintos grados de urgencias y sea capaz de proporcionar
una ayuda a los profesionales en caso necesario


	% Estado del arte
	% 	1. Crítica al estado del arte
	% 	2. Propuesta

	\chapter{Estado del arte}

El software libre y sus licencias \cite{gplv3} ha permitido llevar a cabo una expansión del
aprendizaje de la informática sin precedentes.


	\chapter{Metodología}

Para comenzar el desarrollo de un proyecto primero necesitamos una planificación inicial de la que partir.

\section{Planificación}

La planificación que se va a utilizar está basada en los {\href{https://agilemanifesto.org/iso/es/principles.html}{principios del manifiesto ágil}}
Buscamos un tipo de planificación que se centre en el usuario, para que el propio usuario nos aporte retroalimentación.
Una planificación ágil nos permite adaptarnos al cambio de requisitos haciendo uso de las iteraciones y 
evitando caer en una documentación excesiva llena de diagramas y requisitos innecesarios.
El desarrollo ágil construye el sistema por medio de aportaciones frecuentes de código que aportan valor 
al usuario, es necesario organizar las funcionalidades en bloques y que los cambios siempre surgen de una 
necesidad del usuario.
Para cualquier tipo de planificación es indispensable una herramienta de control de versiones, para ello voy a emplear  git
y GitHub como repositorio para hacer un seguimiento de todas las tareas.
He de mencionar que todos los commits deben de hacer referencia a las tareas preestablecidas, por lo tanto,
podremos saber cuando leamos el código hacía que historia de usuario o issue se refiere
y poder observar el desarrollo en conjunto del proyecto 

\section{Hitos}
Los hitos pueden definirse como metas a las que debemos de llegar durante el desarrollo del proyecto.
Los hitos son una buena forma de limitar los bloques de trabajo del desarrollo y su principal intención es que 
al final de cada hito consigamos un producto mínimamente viable (\textit{PMV}) para que sea independiente e iterativo 
Los hitos que se deben de cumplir están reflejados en el GitHub y son los siguientes: 

\subsection*{\href{https://github.com/RubenDelgadoPareja/TFG-Triage-Inteligente-Consulta-Medica/milestone/1}{Hito 0: Infraestructura inicial y documentación}}

El primer PMV consiste en alcanzar un repositorio inicial donde comenzar a forjar toda la arquitectura del proyecto con una planificación inicial.
Para aceptar el PMV, el repositorio debe de contener:

\begin{itemize}
    \item{User Journeys.}
    \item{Hitos.}
    \item{Historias de Usuario.}
    \item{Comprobadores de gramática y ortografía.}
    \item{Comprobadores de compilación de LaTeX para la documentación.}
    \item{Documentar las secciones de Motivación, Objetivos y Planificación}
\end{itemize}

\subsection*{\href{https://github.com/RubenDelgadoPareja/TFG-Triage-Inteligente-Consulta-Medica/milestone/7}{Hito 1: Modelado del Problema con DDD}}

Este PMV consiste en modular el dominio del problema para facilitar el diseño e implementación.
Para aceptar el PMV la memoria debe de recoger una sección específica sobre cómo se aplicará el DDD la cual incluya: 

\begin{itemize}
    \item {Definición del dominio}
    \item {Definir el lenguaje ubicuo}
    \item {Identificar Entities, Value Objects, Aggregates}
    \item {Definir los Repositories y Services necesarios}
\end{itemize}


\subsection*{\href{https://github.com/RubenDelgadoPareja/TFG-Triage-Inteligente-Consulta-Medica/milestone/2}{Hito 2: Arquitectura básica del sistema de colas de citas}}

Este PMV consiste en diseñar las clases y funciones necesarias para crear citas, almacenarlas en colas y priorizarlas mediante una heurística.

Los criterios de aceptación de este PMV son los tests de cada clase y sus funciones, además de reflejar en la memoria las decisiones tomadas

\section{Historias de Usuarios}
Las historias de usuario son la unidad mínima de funcionalidad que el usuario necesita del sistema, es decir, 
expresa lo que quiere el usuario del sistema por esto lo hace tan importante, se centra completamente en el 
usuario. Encontraremos las Historias de Usuario repartidas por los Hitos del desarrollo del software .

\noindent{Se han definido las siguientes historias de usuario:}

\subsection*{\href{https://github.com/RubenDelgadoPareja/TFG-Triage-Inteligente-Consulta-Medica/issues/19}{[HU-01] Triage Inteligente - Priorizar una cita.}}

Como paciente quiero tener prioridad en el orden de citas médicas.

\subsection*{\href{https://github.com/RubenDelgadoPareja/TFG-Triage-Inteligente-Consulta-Medica/issues/5}{[HU-02] Triage Inteligente - Establecer riesgo.}}

Como sanitario quiero asignarle una prioridad o riesgo a un paciente

\section{Issues}
Las Issues son simples tareas que hay que realizar para completar los Hitos. 
Una historia de Usuario podría ser un conjunto de Issues, aunque en GitHub las he englobado en una sola Issue
Todas las Issues que solucionaré durante el desarrollo para construir el sistema están descritas en GitHub 


\section{Control de calidad}




	% Desarrollo bajo sprints:
	% 	1. Permitir registros y login de usuarios
	% 	2. Desarrollo del sistema de incidencias
	% 	3. Desarrollo del sistema de denuncias administrativas y accidentes
	% 	4. Desarrollo del sistema de croquis
	%   5. Instalación de la aplicación de manera automática
	\UseRawInputEncoding
\chapter{Implementación}

A lo largo de este capítulo se pretende reflejar las herramientas usadas durante todo el proceso de desarrollo y
se argumentará por qué se han tomado estas decisiones, al igual que en el capítulo de \ref{ch:metodología} explicamos
la planificación y forma de trabajo.

La implementación del software se ha dividido en hitos. Cada hito representa un producto mínimamente viable que se publicará.
Los \href{https://github.com/RubenDelgadoPareja/TFG-Triage-Inteligente-Consulta-Medica/milestones}{hitos} han sido definidos en GitHub y explicados en la sección \ref{sec:hitos}.
Cada uno de ellos contiene un grupo de \href{https://github.com/RubenDelgadoPareja/TFG-Triage-Inteligente-Consulta-Medica/issues}{\textit{issues}} que se corresponden
con los distintos problemas que han ido surgiendo durante el desarrollado, cómo los hemos solucionado, qué herramientas se han utilizado y por qué.

\section{Hito 0: Infraestructura inicial y documentación}

El primer PMV consiste en alcanzar un repositorio inicial donde comenzar a forjar toda la arquitectura del proyecto con una planificación inicial.
En el capítulo de metodología ya definimos el PMV a alcanzar para este hito \ref{sb:hito0}.

\subsection{Entorno de Desarrollo Integrado}
El primer problema es encontrar un lugar donde comenzar el desarrollo del proyecto y configuraciones de herramientas, para ello necesitamos un entorno de desarrollo integrado (\textit{IDE}).
He decidido elegir Visual Studio Code para el desarrollo del proyecto debido a que tengo experiencia usándola y me siento cómodo desarrollando en este {\textit{IDE}}.
Otro motivo fundamental es que es gratuita, además contiene una gran cantidad de extensiones y configuraciones lo que facilita programar, también es ligero y multiplataforma.
Me parece muy cómodo y versátil frente a otras opciones.

\subsection{Memoria}
La memoria es un documento muy importante en el desarrollo de un proyecto, ya que, es el reflejo de todo el trabajo realizado y es necesario para la evaluación del proyecto.
Debemos de encontrar una herramienta que nos permita escribir la memoria de forma cómoda y eficiente. Además de comprobar constantemente que lo escrito esté correcto.

\paragraph*{Compilación de LaTeX}
La memoria se ha escrito empleando el lenguaje LaTeX, que es un sistema de composición de documentos bastante empleado
en el ámbito académico y científico, es necesario para poder generar la documentación final, agrupando cada capítulo, sección y
gráfica. Para generar este archivo final debemos ejecutar el comando \textit{all} del Makefile

\paragraph*{Corrector ortográfico gramatical}
Para verificar activamente que lo escrito en la memoria esté correcto se ha empleado una herramienta de LaTeX llamada \href{https://github.com/sylvainhalle/textidote}{TeXtidote}.
Esta herramienta nos permite corregir la ortografía, gramática y semántica de la documentación escrita. Se puede ejecutar la corrección a través del comando \textit{check} del Makefile,
generando un borrador con la documentación marcando los errores y posibles cambios. Además, se ha añadido este comando a los hooks de GitHub para ejecutarlo antes de cada commit.

\subsection{Código}
El código es imprescindible para el desarrollo, debemos encontrar un lenguaje que satisfaga las necesidades de la \ref{sc:planificación}.
Además, es necesario tener un constante análisis estático para mantener un código limpio, coherente y libre de errores.
Finalmente, para aportar valor al producto he decidido automatizar estos procesos con la integración continua, mencionado en la sección \ref{sc:control-calidad}.

\subsubsection*{Lenguaje de Programación}
Para el desarrollo de un proyecto es esencial decidir con qué herramienta principal, en este caso el lenguaje, se va a crear. De entre mi pila de lenguajes de programación he
escogido Typescript junto con Node.js, ya que, estoy trabajando con él. Creo que da una gran robustez gracias al tipado y apoya mucho al Diseño Dirigido al Dominio, no como otros lenguajes como Java o Ruby
Cuenta con una gran batería de librerías que agiliza el desarrollo y permite la implementación rápida de características.

\subsubsection*{Análisis estático}
Para garantizar la calidad, la eficiencia y sostenibilidad del código a largo plazo, necesitamos un análisis estático del código.
Se ha decidido elegir ESLint con Typescript, ya que es un \textit{linter} que se emplea para el lenguaje de JavaScript, pero se puede aprovechar con TypeScript.
ESLint te permite aplicar reglas personalizadas, reglas predefinidas y configuraciones de estilo para garantizar que tu código sea consistente y libre de problemas.
No se ha escogido TSLint, por el hecho de que, está obsoleta y no tiene mantenimiento, también se ha barajado SonarQube y CodeClimate.

\subsubsection*{Integración Continua}
La integración continua es una práctica de desarrollo de software que tiene como objetivo principal mejorar la calidad del código, aumentar la eficiencia y agilizar el proceso de entrega de software.
Las diferentes opciones que encontramos para el cliente o local pueden ser git, SVN, Mercurial y para el servidor donde se harán las pruebas y la integración se presentan las siguientes
alternativas Jenkins, Travis, Circle, GitHub Actions. He decido escoger git para el cliente y GitHub Actions para el servidor porque me resultan fáciles de configurar y escalabilidad.
Gracias a esta metodología podemos automatizar la ejecución de test, linter y corrector ortográfico, de esta manera comprobaremos antes de ciertas acciones con los hooks de GitHub que todo está
correctamente. También se había barajado la opción de utilizar un paquete llamado \textit{husky} para Node.js, pero no se ha escogido finalmente para aprender nuevas tecnologías.

\section{Hito 1: Modelado del Problema con DDD}

En este segundo PVM nos adentramos en el modelado del problema descrito en la sección \ref{sc:hito1}.
Pero antes debemos investigar sobre el Diseño Dirigido al Dominio y aplicarlo correctamente al dominio del problema. Además, con el desarrollo del
código surge la necesidad de realizar test unitarios e integración, mencionado en la sección \ref{sc:control-calidad}.

\subsection{Diseño}
La implementación del código estará guiada por el Diseño Dirigido al Dominio, esto significa que el diseño del código se basará en el dominio del problema.
Podemos encontrar una pequeña introducción al DDD en la sección \ref{sc:diseño}, no obstante es necesario explicar cómo se ha aplicado al proyecto.

Inicialmente, se deben localizar los subdominios y los contextos acotados. La arquitectura del proyecto se ha dividido en módulos que representan los distintos contextos acotados.
Cada módulo tendrá su propio conjunto de clases que representan esa parte del dominio y debe de tener sentido por sí mismo. Los distintos módulos también podrán comunicarse entre ellos.
Inicialmente, se van a crear lo mínimo necesario para el cliente, después poco a poco se irá aportando valor al producto, añadiendo más funcionalidades y mejoras.

\subsection{Test unitarios e integración}
Para garantizar la calidad del código, es necesario realizar test unitarios e integración. Los test unitarios son pruebas que se ejecutan a una unidad de código, por ejemplo, una función o un método.
Las herramientas que se han considerado para testear han sido Cypress, Mocha y Jest.
En general, todas son buenas opciones, pero se ha escogido Jest por su facilidad de uso, su integración con TypeScript y su gran comunidad.


	% Presupuesto

	% Conclusiones
	\UseRawInputEncoding
\chapter{Conclusiones y trabajos futuros}



	% Trabajos futuros



	\newpage
	\bibliography{bibliografia}
	\bibliographystyle{plain}

\end{document}

